\documentclass[../main.tex]{subfiles}
\graphicspath{{\subfix{../image/}}}
\begin{document}
\chapter{积分曲线与流}\label{chapter:ch09-积分曲线与流}

本章继续研究向量场。与光滑向量场相关的主要几何对象是其积分曲线，即每一点处的速度都等于该点处向量场值的光滑曲线。流形上给定向量场的所有积分曲线的集合确定了该流形（开子集）的一族微分同胚，称为\textbf{流}。例如，任意光滑的 $\mathbb{R}$-作用都是流；但也存在不是 $\mathbb{R}$-作用的流，因为微分同胚不一定对每个 $t\in\mathbb{R}$ 都在整个流形上有定义。

本章的主定理，即流的基本定理，断言每个光滑向量场确定唯一的、从每一点出发的极大积分曲线，而所有此类积分曲线的集合确定唯一的极大流。其证明是常微分方程解的存在性、唯一性和光滑性定理（见附录 D）的应用。

证明基本定理后，我们探讨向量场和流的一些性质。首先研究向量场生成全局流的条件。然后说明如何沿向量场从初始子流形"流出"来创建较大子流形的有用参数化。接下来考察流的局部行为，发现在向量场为零的点（对应于流的平衡点）处的行为与在其非零点处的行为非常不同，在后者处流局部看起来像沿平行坐标直线的平移。

然后引入 Lie 导数，这是一种坐标无关地计算一个向量场沿另一个向量场的流的变化率的方式。它导出向量场、其 Lie 括号及其流之间的深刻联系。特别地，我们将证明两个向量场有交换的流，当且仅当它们的 Lie 括号为零。基于这一事实，可以证明光滑局部标架可表示为坐标标架的充要条件。然后讨论本章的一些结果如何推广到流形上的含时向量场。

本章最后一节描述流对一阶偏微分方程研究的重要应用。

\subfile{Chapter-09/section-01.tex}
\subfile{Chapter-09/section-02.tex}
\subfile{Chapter-09/section-03.tex}
\subfile{Chapter-09/section-04.tex}
\subfile{Chapter-09/section-05.tex}
\subfile{Chapter-09/section-06.tex}
\subfile{Chapter-09/section-07.tex}
\subfile{Chapter-09/section-08.tex}
\subfile{Chapter-09/section-09.tex}
\end{document}